% Generated draft: requires mathematical and source review.
\documentclass[11pt]{article}
\usepackage[margin=25mm]{geometry}
\usepackage{fontspec}
\setmainfont{Latin Modern Roman}
\usepackage{amsmath,amssymb,mathtools,unicode-math}
\setmathfont{Latin Modern Math}
\usepackage{enumitem,needspace,xurl,hyperref}
\setlength{\parindent}{0pt}
\setlength{\parskip}{5pt}
\setlength{\emergencystretch}{3em}
\newcommand{\N}{\mathbb{N}}
\newcommand{\Z}{\mathbb{Z}}
\newcommand{\Q}{\mathbb{Q}}
\newcommand{\R}{\mathbb{R}}
\newcommand{\C}{\mathbb{C}}
\newcommand{\F}{\mathbb{F}}
\newcommand{\A}{\mathbb{A}}
\newcommand{\PP}{\mathbb{P}}
\newcommand{\E}{\mathbb{E}}
\newcommand{\eps}{\varepsilon}
\newcommand{\dd}{\,\mathrm{d}}
\begin{document}
\section*{UnsolvedMath ID: 30000245}
\textbf{Source problem code:} catalogue entry
\subsection{Definitions and mathematical statement}
For a knot $K$, let $H_N^{i,j}(K)$ and $\overline H_N^{i,j}(K)$ denote its reduced and unreduced $\mathfrak{sl}(N)$ Khovanov--Rozansky homologies, with the source's coefficient and grading conventions. Write their Poincar\'e polynomials as
\[
 \operatorname{KhR}_N(K;q,t)=\sum_{i,j}\dim H_N^{i,j}(K)\,q^jt^i,
 \qquad
 \overline{\operatorname{KhR}}_N(K;q,t)=
                      \sum_{i,j}\dim\overline H_N^{i,j}(K)\,q^jt^i.
\]
The precise conjecture is that for every knot $K$ there exist finite Laurent polynomials
\[
 \overline{\mathcal P}_K\in\Z[a^{\pm1},q^{\pm1},t^{\pm1}],\qquad
 \mathcal P_K\in\Z_{\ge0}[a^{\pm1},q^{\pm1},t^{\pm1}],
\]
and an integer $N_0(K)$ such that for every integer $N\ge\max\{2,N_0(K)\}$,
\[
 (q-q^{-1})\overline{\operatorname{KhR}}_N(K;q,t)
       =\overline{\mathcal P}_K(q^N,q,t),\qquad
 \operatorname{KhR}_N(K;q,t)=\mathcal P_K(q^N,q,t).
\]
The first identity is equivalently the source's unreduced quotient by $q-q^{-1}$. These are formal identities in the grading variables. The two universal polynomials depend on $K$ but not on $N$, and each has only finitely many nonzero terms; the stabilization threshold may depend on $K$. The reduced coefficients must be nonnegative, while the unreduced numerator is allowed signed integer coefficients. No nonnegativity condition is added to that numerator, nor is the assertion extended to all small $N$.

The extract does not spell out homology coefficient field, reduced-basepoint convention or grading shifts. They must be fixed consistently with the displayed source normalizations. The finite unreduced numerator should not be identified without justification with the Poincar\'e polynomial of a finite-dimensional universal unreduced triply graded homology.

The literature assessment checked on 21 August 2026 reports important universal triply graded and Khovanov--Rozansky constructions and a 2022 proof of the predicted spectral sequence, but does not identify that theorem with the full finite-polynomial statement above for every knot. Its partial-progress claim and the distinction between spectral-sequence existence and these exact stable specializations are retained as supplied dated evidence, requiring independent verification.
\subsection{Short English statement}
For each knot, find two finite universal Laurent polynomials giving reduced and unreduced $\mathfrak{sl}(N)$ homology polynomials for all sufficiently large $N$. Only the reduced universal polynomial is required to have nonnegative coefficients, and the unreduced specialization includes the stated denominator.
\subsection{Sources}
\begin{itemize}
\item \emph{Geometric Topology and Connections with Quantum Field Theory} (2006), Oberwolfach Reports, DOI 10.4171/owr/2005/27, pp. 1505--1546, \url{https://doi.org/10.4171/owr/2005/27}; proposed year supplied as 2005.
\item Literature assessment checked 21 August 2026: Dunfield--Gukov--Rasmussen, \url{https://arxiv.org/abs/math/0505662}; ``Proof of DGR conjecture,'' \url{https://arxiv.org/abs/2210.00878}; ``Positive torus-knot formula,'' \url{https://msp.org/gt/2022/26-1/gt-v26-n1-p02-p.pdf}.
\item Local UnsolvedMath registry, permanent ID $30000245$, exact source code \texttt{catalogue entry}, created 25 August 2026 and updated 16 September 2026. Supplied verification restores both finite-polynomial assertions, their distinct coefficient rings and the large-$N$ specializations.
\end{itemize}
\end{document}
